\documentclass[11pt]{article}

% --------------------------------------------------
% LuaLaTeX setup
% --------------------------------------------------
\usepackage{fontspec}
\setmainfont{Latin Modern Roman}

\usepackage{geometry}
\usepackage{microtype}
\usepackage{setspace}
\usepackage{amsmath,amssymb,amsthm}
\usepackage{csquotes}
\usepackage{hyperref}

\geometry{margin=1in}
\setstretch{1.15}

\newtheorem{definition}{Definition}
\newtheorem{proposition}{Proposition}

% --------------------------------------------------
% Title
% --------------------------------------------------
\title{Constraints Without Commitment:\\
On the Limits of Algebraic Unification in Contemporary Deep Learning}
\author{Flyxion}
\date{December 2025}

\begin{document}
\maketitle

% ==================================================
\section{Introduction}
% ==================================================

Recent work in Categorical Deep Learning (CDL) proposes a unifying mathematical
language for neural network design, capable of bridging the long-standing divide
between top-down constraints and bottom-up implementation. By appealing to
category theory—specifically to universal algebra in a 2-category of parametric
maps—this framework offers a principled account of architectural invariance,
recursion, and weight tying.

The ambition of this project is substantial and, in many respects, successful.
Practices that previously appeared as engineering heuristics are reinterpreted
as algebraic necessities. Entire classes of models—geometric, recurrent, and
message-passing—are shown to instantiate a shared compositional structure.

This essay does not contest these achievements. Instead, it argues that the very
success of categorical unification reveals a boundary. While CDL excels at
formalizing \emph{constraints} on computation, it does not formalize
\emph{commitment} in the sense required for agency, history, and accountability.
The distinction is not cosmetic. It marks the limit of what can be expressed in
purely algebraic terms.

The argument proceeds by strengthening the technical foundation of this claim,
deepening its philosophical implications, and engaging constructively with the
framework’s genuine contributions.

% ==================================================
\section{Accountable Lawfulness}
% ==================================================

The preceding analysis suggests a precise characterization of the systems
described by categorical deep learning. They are maximally lawful and minimally
accountable. They perfectly respect algebraic priors, internalize structure, and
generalize efficiently. At the same time, they lack the internal resources to
record refusal, bear loss, or bind themselves to history.

Non-invertible computation allows such systems to destroy information, but not to
incur obligation. Actions change representations, but they do not settle futures.
The result is a form of intelligence that participates in computation but not in
history.

This diagnosis does not undermine the value of categorical unification. Rather,
it clarifies its scope. The framework explains how systems act. It does not
explain why they never give, refuse, or stop. That gap is not accidental; it is
the mark of a boundary between lawful structure and accountable agency.


% ==================================================
\section{What Categorical Deep Learning Gets Right}
% ==================================================

Before turning to critique, it is essential to acknowledge the genuine advances
offered by the categorical approach.

First, CDL achieves a real unification. Prior to this work, geometric deep
learning, recursive architectures, and weight sharing were treated as distinct
design paradigms. CDL demonstrates that these are instances of a common algebraic
structure: structure-preserving maps between parameterized computational spaces.
This clarification is not merely aesthetic; it explains why certain designs
generalize well and others fail.

Second, the framework explains the remarkable data efficiency of models with
built-in invariances. By encoding permutation equivariance or other symmetries
directly into the architecture, the hypothesis space is dramatically reduced.
Empirically, this yields near-exponential reductions in sample complexity—a
result that is both theoretically satisfying and practically significant.

Third, the identification of the \emph{carry problem} in graph neural networks
represents a nontrivial technical insight. The difficulty GNNs have with
discrete, cascading transitions is not obvious from standard formulations. The
connection to nontrivial fibrations in topology underscores the mathematical
depth of the observation.

These achievements justify taking the framework seriously on its own terms.

% ==================================================
\section{Algebraic Constraints and Their Scope}
% ==================================================

To understand the limits of categorical unification, it is necessary to examine
the progression of algebraic structures it employs.

\subsection{From Groups to Categories}

\begin{definition}
A \emph{group} is a set $G$ equipped with a binary operation $\cdot$ satisfying
closure, associativity, the existence of an identity element, and invertibility.
\end{definition}

Group actions underlie much of geometric deep learning. Rotations, reflections,
and permutations act on inputs in ways that can always be undone. No information
is lost.

\begin{definition}
A \emph{monoid} is a set $M$ with an associative binary operation and an identity,
but without the requirement of invertibility.
\end{definition}

Monoids naturally model computation. Algorithms such as shortest-path solvers map
many distinct inputs to a single output, destroying information irreversibly.

\begin{definition}
A \emph{category} consists of objects, morphisms between objects, identity
morphisms, and associative composition, without requiring that all morphisms be
composable or invertible.
\end{definition}

Categories generalize both structures, allowing heterogeneous composition and
non-invertible processes.

\subsection{Quotients and Erasure}

Consider Dijkstra’s algorithm. Distinct weighted graphs $G_1, G_2, \dots$ may all
map to the same shortest-path tree. Formally, the algorithm defines a quotient
map from the space of graphs to equivalence classes determined by path length.
The categorical formalism captures this quotienting cleanly.

What it does not capture is which representative was chosen. The computation is
lawful, but its result erases history. This distinction—between preserving
structure and preserving commitment—will recur.

\subsection{Lawfulness vs. Commitment}

Structure-preserving maps define \emph{lawful computation}. Commitment, by
contrast, involves selection from equivalence classes with consequences for
future action. Algebraic invariance can describe the former, but it is silent on
the latter.

% ==================================================
\section{Constraints as Homomorphism Requirements}
% ==================================================

Within the categorical framework, a constraint is not a numerical restriction on
parameters but a semantic requirement: a neural layer must act as a homomorphism
between algebras. This reframes architectural design as the problem of selecting
endo-functors whose algebras encode the intended computational behavior.

In this sense, constraints remain meaningful even under reparameterization.
Two-morphisms in the 2-category of parametric maps permit parameters to vary, but
only insofar as the induced morphism preserves the algebraic structure. A
reparameterization that violates the homomorphism condition is not merely
suboptimal; it is ill-typed within the theory.

This interpretation clarifies why categorical deep learning distinguishes
between syntactic freedom and semantic rigidity. Syntax may vary freely, but
semantics—captured by the algebra—must remain invariant. The framework therefore
addresses the question of how computation is shaped, not how future action is
bound.

The distinction is decisive. Homomorphism requirements constrain what a system
\emph{can compute}, but they do not constrain what a system \emph{must stand
behind}. Constraint, in this sense, is alethic rather than deontic: it governs
possibility, not obligation.


% ==================================================
\section{Recursion, Folding, and the Erasure of History}
% ==================================================

Recursion is modeled categorically as a fold: a homomorphism from an inductively
defined data structure to a result. In this view, the history of a computation
is preserved only insofar as it is encoded in the structure of the input. The
execution trace itself is erased in favor of the final value.

Similarly, the unrolling of recurrent networks via copy maps encodes repetition
rather than lineage. Weight tying ensures that the same computation is applied
at each step, but it does not record the fact that a particular path was taken
or refused. The endofunctor governing the recurrence specifies lawful evolution,
not historical memory.

What survives is structure; what disappears is biography.

% ==================================================
\section{Internalized Reasoning and the Absence of Obligation}
% ==================================================

A central claim of the categorical approach is that models should internalize
algorithms rather than rely on external tools. Internalization yields stability,
efficiency, and correctness guarantees. Carrying in addition or conservation
laws in physics should be reflected in the model’s internal geometry, not
recovered through repeated querying.

However, what is internalized are procedures, not consequences. Cost appears as
a computational budget, not as an obligation. Correctness is mathematical, not
normative. The framework provides no mechanism by which a system could absorb
risk, incur debt, or be worse off as a result of its actions.

Internalized reasoning remains non-dissipative. 

% ==================================================
\section{Internalization, Stability, and the Limits of Consequence}
% ==================================================

The categorical program emphasizes the internalization of algorithms as a source
of stability and efficiency. Internalized reasoning avoids the brittleness of
external tool calls, which require accurate input prediction and repeated
re-invocation. In this sense, internalization reduces computational cost and
improves robustness.

However, what is internalized is the procedure, not the consequence. Stability
is achieved by aligning neural machinery with symbolic computation, not by
encoding responsibility for failure. Correctness guarantees remain mathematical,
expressed as post-conditions or invariants, rather than normative commitments.

The framework acknowledges uncertainty by suggesting that systems might be
trained to back off when computational demands exceed their capacity. Yet such
backing off remains instrumental. It is triggered by resource limits, not by a
standing obligation to refuse.

Internalization thus stabilizes reasoning without internalizing stake. The system
becomes better at acting correctly, not at being bound by having acted.


% ==================================================
\section{The Carry Problem and Emergent Obligation}
% ==================================================

The carry operation in arithmetic provides a concrete illustration of the
distinction between local structure and global consequence.

Adding digits modulo 10 treats each digit independently. Ordinary addition,
however, is not independent: a transition from $9 \rightarrow 0$ in one digit
forces an increment in the next. The obligation is not in the state, but in the
transition.

Graph neural networks struggle with such phenomena because message passing
captures local state changes but not cascading obligations. The Hopf fibration
$S^3 \rightarrow S^2$, with fiber $S^1$, provides a continuous analogue: the
total space is not decomposable into independent components.

Once a carry occurs, it cannot be undone without altering the entire structure.
The categorical framework captures the algebraic structure of carrying, but it
does not record the fact that a carry has occurred as a binding historical event.

% ==================================================
\section{Parametric Maps and Asymmetric Loss}
% ==================================================

The 2-category of parametric maps, Para, makes parameters explicit objects of
study. Objects are pairs $(A, P)$, where $A$ is a space and $P$ a parameter
space. Morphisms are parameterized functions $f : P \times A \rightarrow B$.
Two-morphisms correspond to reparameterizations preserving structure.

Weight tying arises naturally via copy maps. This symmetry explains why shared
parameters induce coordinated behavior.

However, this symmetry also reveals a limitation. There is no notion of ownership
or unilateral commitment. All admissible reparameterizations that preserve
structure are equally valid. Loss functions update parameters symmetrically.

In game theory, commitment involves irrevocably constraining one’s own future
actions to change others’ incentives. This requires burned bridges: regions of
parameter space that cannot be re-entered. Para contains no such mechanism.

Loss exists as information destruction, but no agent bears it.

% ==================================================
\section{Symmetry of Parameters and the Impossibility of Ownership}
% ==================================================

The 2-category of parametric maps makes parameters explicit but treats them as
neutral mathematical objects. Morphisms relate parameterized spaces, and
two-morphisms reparameterize them while preserving structure. At no point does
the framework assign control, ownership, or agency to a particular party.

Asymmetry arises only at the level of computation: non-invertible morphisms
destroy information, and certain natural transformations exist in one direction
but not the other. This asymmetry, however, is internal to the diagram. It does
not imply that any agent bears unilateral loss.

Loss is therefore endogenous to the morphism rather than exogenous to an agent.
Data may be discarded, but no one is worse off. Unilateral commitment—understood
as the irrevocable restriction of one’s own future actions—is not representable
within this symmetric parameter space.

This explains why cooperation in the framework is purely compositional. Joint
behavior is correct when it is isomorphic to the composition of parts. What is
missing is the possibility of binding oneself so that others may rely on that
binding.


% ==================================================
\section{Refusal as a Missing Primitive}
% ==================================================

Optimization frameworks permit weak negation: actions are avoided because they
are costly. Such negation can always be overridden with sufficient resources.

Refusal, by contrast, is strong negation: an action is forbidden regardless of
utility. This distinction maps onto the philosophical divide between alethic
modality (what is possible) and deontic modality (what is permitted).

Category theory is fundamentally alethic. It describes what transformations are
lawful. Commitment and refusal are deontic; they persist independently of
possibility.

Refusal also requires history. Having refused an action at time $t_1$ must
constrain options at $t_2$. Absent a state that records refusal, the system
remains perpetually revisable.

Purely structural models thus inhabit a form of permanent bad faith: past
choices never bind the present.

% ==================================================
\section{Weak and Strong Negation}
% ==================================================

The framework permits a form of weak negation. A system may decline to answer a
question because it exceeds a computational budget or violates a learned prior.
Such negation is contingent and revisable. With sufficient resources or a change
in context, the refusal can be overridden.

Strong negation is categorically different. To refuse in the strong sense is to
forbid an action regardless of cost. Such refusal persists across contexts and
binds future behavior. It is deontic rather than alethic.

Categorical constraints are necessarily weak in this sense. They describe what
cannot occur given the structure of the computation, not what must not occur
despite being possible. Treating refusal as first-class would therefore require
a state that records negation as an event, not merely as the absence of a
morphism.

Introducing such states would break the assumption that computation always
composes. Some morphisms would lead to terminal conditions that are not identities
but prohibitions.

% ==================================================
\section{Time, Reversibility, and Worldhood}
% ==================================================

Categories naturally model processes: continuous transformations composed over
time. Events, however, are breaks—partitions of time into before and after.

Optimization lacks terminality. Local minima are not ends; they are pauses.
Games, by contrast, have terminal states where no further moves are legal.

Non-invertible morphisms capture logical irreversibility, but not temporal
irreversibility. Two identical states reached via different histories are
structurally indistinguishable, yet historically distinct.

Worldhood requires stake. A world is not merely a state space, but a space of
concern in which actions matter. Structure without stake yields process without
world.

% ==================================================
\section{Cooperation, Commitment, and Common Knowledge}
% ==================================================

Shared parameters enable structural coordination, but not trust. If agents can
secretly reparameterize, cooperation collapses.

Game-theoretic commitment requires common knowledge of irreversible moves.
Agreements without enforcement are, as Hobbes observed, merely words.

Categorical alignment provides the grammar of cooperation, but not its binding.
This explains why external tools and verifiers remain necessary: the system
cannot commit to itself.

% ==================================================
\section{Accountable Lawfulness}
% ==================================================

The preceding analysis suggests a precise characterization of the systems
described by categorical deep learning. They are maximally lawful and minimally
accountable. They perfectly respect algebraic priors, internalize structure, and
generalize efficiently. At the same time, they lack the internal resources to
record refusal, bear loss, or bind themselves to history.

Non-invertible computation allows such systems to destroy information, but not to
incur obligation. Actions change representations, but they do not settle futures.
The result is a form of intelligence that participates in computation but not in
history.

This diagnosis does not undermine the value of categorical unification. Rather,
it clarifies its scope. The framework explains how systems act. It does not
explain why they never give, refuse, or stop. That gap is not accidental; it is
the mark of a boundary between lawful structure and accountable agency.

% ==================================================
\section{Analysis}
% ==================================================

Categorical Deep Learning succeeds at unifying constraints and implementation. It
formalizes intelligence as lawful structure and explains why certain architectures
generalize so effectively.

What it does not formalize is commitment. Constraints shape computation, but they
do not bind future action. Irreversibility appears as information loss, not
obligation. Refusal remains instrumental rather than principled.

If computation explains how systems act, it does not explain why they never give.
Crossing that boundary requires a theory in which actions can settle, costs can
be borne, and history can accumulate.

% ==================================================
\section{From Lawful Computation to Eventful Structure}
% ==================================================

The preceding analysis has shown that Categorical Deep Learning provides a
complete account of lawful computation: it unifies constraints and
implementations by identifying architectures with structure-preserving maps.
Within this framework, invariance, recursion, weight tying, and non-invertible
computation are explained as algebraic necessities rather than engineering
choices.

At the same time, this completeness marks a boundary. The categorical formalism
is intentionally indifferent to how a computation is realized in time, to which
events occurred, and to whether alternatives were ever explicitly refused.
Reparameterization preserves lawfulness precisely because the framework is
designed to quotient away implementation history.

What is required to move beyond this boundary is not a revision of the algebraic
account, but a change in level. Rather than enriching architectures directly,
one can descend to a description in which implementations are treated as
eventful processes, and only subsequently recover algebraic structure by
abstraction. In such a setting, histories are primary rather than incidental,
and irreversibility appears as an explicit feature of state rather than as
information loss alone.

This motivates the introduction of a field-theoretic intermediate description
in which scalar accumulations, vector flows, and entropy production summarize
the thermodynamic consequences of event histories. From this level, both
geometric constraints and categorical laws can be recovered as distinct
abstractions that forget different aspects of the underlying dynamics.

The following sections make this claim precise. They show how an event-sourced
implementation calculus gives rise to an RSVP-style field theory, and how
Geometric Deep Learning and Categorical Deep Learning emerge as orthogonal
quotients of that field-level description. In this way, accountability and
commitment are not appended to categorical structure, but shown to arise from
what categorical abstraction necessarily leaves behind.


% ==================================================
\section{From Eventful Implementations to Algebraic Architectures}
% ==================================================

The categorical deep learning framework is often presented as a bridge between
constraints and implementation. In this work, however, the direction of
derivation is reversed. Rather than treating categorical structure as
foundational and implementations as instances, we begin from an
event-sourced implementation calculus and show how both Geometric Deep Learning
and Categorical Deep Learning arise as abstractions of a more primitive
field-theoretic description.

At the base level lies an operational calculus in which structures are generated
by irreversible events. Objects are not primitive sets or vector spaces but
stabilized outcomes of event histories. Composition is temporal rather than
algebraic, and identity corresponds to the absence of events rather than a
neutral morphism. This level describes implementation in the strongest sense:
what exists is only what has been constructed.

From this eventful substrate, a field-theoretic description emerges by
coarse-graining over fine-grained event order. Distinct histories are identified
when they induce the same scalar accumulations, vector flows, and entropy
production. The result is a relativistic scalar–vector–entropy field theory in
which dynamics are governed by conservation, flow, and dissipation rather than
by discrete computational steps. At this level, commitment appears as
irreversibility in the entropy field, and history is preserved only through its
thermodynamic consequences.

Geometric Deep Learning is obtained by a further abstraction that forgets
thermodynamic content while retaining symmetry structure. When RSVP field
configurations are quotiented by group actions—permutations, rotations, or more
general equivariances—the remaining structure specifies which transformations
leave the dynamics invariant. Message passing, equivariant layers, and symmetry-
respecting architectures arise as implementations that are compatible with this
symmetry-reduced description. In this sense, geometric constraints are not
primitive assumptions but residual structure left after discarding energetic
and entropic detail.

Categorical Deep Learning arises by a different abstraction of the same RSVP
substrate. Instead of quotienting by symmetry, one quotients by lawful
replayability of computation. Event histories are identified when they implement
the same abstract computation, regardless of parameterization or concrete
realization. What survives this identification are algebras over endofunctors,
parametric maps, and structure-preserving transformations between them. Weight
tying, recursion, and non-invertible computation appear here as algebraic
properties of stable implementations, not as foundational axioms.

On this view, Geometric Deep Learning and Categorical Deep Learning are not
competing foundations but orthogonal projections of a common underlying theory.
Both are derived from RSVP by forgetting different aspects of the eventful
substrate: geometry forgets thermodynamic commitment, while category theory
forgets concrete event history. Neither abstraction recovers the full structure
of the underlying implementation calculus, but each captures a distinct class
of invariants relevant to learning.

The ordering is therefore essential. RSVP is not recovered from categorical deep
learning; categorical deep learning is recovered from RSVP by abstraction. What
categorical deep learning unifies—constraints and implementation—has already
been unified at the level of eventful dynamics. The categorical framework
formalizes the lawful remainder after irreversible, field-level commitments have
been erased.

% ==================================================
\section{Mathematical Derivation as a Chain of Quotients and Forgetful Functors}
% ==================================================

This section makes precise the claim that (i) a field-theoretic description
arises as a coarse-graining of eventful implementations, and that (ii) both
Geometric Deep Learning and Categorical Deep Learning arise as further
abstractions of that field-level description. The construction is expressed as a
sequence of functors that forget different aspects of structure.

% --------------------------------------------------
\subsection{Eventful implementations as a history category}
% --------------------------------------------------

Let $\mathsf{Ev}$ be a fixed set of primitive event types (e.g.\ POP, MERGE, LINK,
COLLAPSE) equipped with admissibility conditions that determine which events may
follow which.

\begin{definition}[Histories and implementations]
A \emph{history} is a finite (or countable) admissible sequence
$h = (e_1,\dots,e_T)$ with $e_t \in \mathsf{Ev}$.
An \emph{implementation state} is an object $x$ produced by replaying a history
from a distinguished initial state $x_0$, written $x = \mathrm{Replay}(h;x_0)$.
\end{definition}

We package histories into a category in the usual ``free category'' manner.

\begin{definition}[History category $\mathsf{Hist}$]
Define a category $\mathsf{Hist}$ whose objects are implementation states and
whose morphisms are histories that transform one state to another:
\[
h : x \to y
\quad \text{iff} \quad
y = \mathrm{Replay}(h;x).
\]
Composition is concatenation of histories (when admissible), and identities are
empty histories.
\end{definition}

Intuitively, $\mathsf{Hist}$ is the most concrete level: it remembers \emph{which
events occurred and in what order}. All later structures are obtained by
forgetting some of this detail.

% --------------------------------------------------
\subsection{Coarse-graining to RSVP: thermodynamic equivalence}
% --------------------------------------------------

Let $\mathsf{Fld}$ denote a category of RSVP field configurations. For present
purposes, an object of $\mathsf{Fld}$ may be taken as a triple
\[
X = (\Phi,\mathbf{v},S)
\]
where $\Phi$ is a scalar field, $\mathbf{v}$ a vector field, and $S$ an entropy
(or irreversibility) field; morphisms in $\mathsf{Fld}$ are admissible field
updates (e.g.\ time-evolution steps) compatible with the chosen dynamics.

\begin{definition}[Thermodynamic coarse-graining functor]
A \emph{thermodynamic coarse-graining} is a functor
\[
\mathcal{Q}_{\mathrm{RSVP}} : \mathsf{Hist} \to \mathsf{Fld}
\]
such that, for any state $x$, the object $\mathcal{Q}_{\mathrm{RSVP}}(x)$ is the
field configuration $(\Phi_x,\mathbf{v}_x,S_x)$ obtained by aggregating the
history-sourced quantities of $x$, and for any history $h:x\to y$,
$\mathcal{Q}_{\mathrm{RSVP}}(h)$ is the induced field update.
\end{definition}

The key property is that $\mathcal{Q}_{\mathrm{RSVP}}$ identifies many distinct
histories as the same field trajectory.

\begin{definition}[RSVP equivalence of histories]
Define $x \sim_{\mathrm{RSVP}} y$ iff
$\mathcal{Q}_{\mathrm{RSVP}}(x) = \mathcal{Q}_{\mathrm{RSVP}}(y)$ in $\mathsf{Fld}$.
Likewise define $h \sim_{\mathrm{RSVP}} h'$ iff
$\mathcal{Q}_{\mathrm{RSVP}}(h) = \mathcal{Q}_{\mathrm{RSVP}}(h')$.
\end{definition}

\begin{proposition}[Field-level description is a quotient of implementations]
$\mathcal{Q}_{\mathrm{RSVP}}$ factors through the quotient of $\mathsf{Hist}$ by
$\sim_{\mathrm{RSVP}}$. In particular, the field description forgets fine-grained
event order and retains only its thermodynamic consequences.
\end{proposition}

\begin{proof}
By definition, objects and morphisms identified by $\sim_{\mathrm{RSVP}}$ have
the same image in $\mathsf{Fld}$. Therefore $\mathcal{Q}_{\mathrm{RSVP}}$ is
constant on equivalence classes and factors through the corresponding quotient
construction.
\end{proof}

This is the formal sense in which RSVP is \emph{derived} from implementation
histories: it is what remains invariant under thermodynamic coarse-graining.

% --------------------------------------------------
\subsection{Deriving Geometric Deep Learning from RSVP by symmetry reduction}
% --------------------------------------------------

Let $G$ be a group acting on the relevant field configurations (e.g.\ permutations
of nodes, rotations, translations), written as
\[
\rho : G \to \mathrm{Aut}_{\mathsf{Fld}}(X), \qquad g \mapsto \rho(g).
\]

\begin{definition}[Symmetry reduction / orbit functor]
Define a functor
\[
\mathcal{Q}_{\mathrm{GDL}} : \mathsf{Fld} \to \mathsf{Fld}/G
\]
sending a configuration $X$ to its orbit $[X]=\{ \rho(g)X : g\in G\}$ and sending
a morphism $f:X\to Y$ to its induced map $[f]:[X]\to [Y]$ whenever this is
well-defined (i.e.\ $f$ is $G$-equivariant).
\end{definition}

Architectures in Geometric Deep Learning may be characterized as maps that
descend to the quotient and therefore respect these symmetries.

\begin{definition}[Equivariant maps as admissible implementations]
A map $f:X\to Y$ is \emph{$G$-equivariant} if
\[
f(\rho(g)X)=\rho(g)f(X)\quad \text{for all }g\in G.
\]
Equivariant maps are precisely those that induce well-defined morphisms
$[f]:[X]\to[Y]$ in $\mathsf{Fld}/G$.
\end{definition}

\begin{proposition}[GDL as a quotient of RSVP]
The symmetry-reduced description $\mathsf{Fld}/G$ forgets thermodynamic and
entropic detail beyond what is required for equivariance, retaining only the
invariance constraints. Equivariant neural layers are implementations of
morphisms in $\mathsf{Fld}/G$.
\end{proposition}

\begin{proof}
Passing from $X$ to $[X]$ identifies all $G$-related configurations; thus only
invariants under $G$ remain. A layer implements a morphism on orbits precisely
when it is equivariant, i.e.\ it respects the quotient.
\end{proof}

Thus Geometric Deep Learning is obtained from RSVP by forgetting more: it discards
field ``substance'' and retains symmetry-encoded constraints.

% --------------------------------------------------
\subsection{Deriving Categorical Deep Learning from RSVP by lawfulness quotients}
% --------------------------------------------------

To extract Categorical Deep Learning, we do \emph{not} quotient by symmetry but by
\emph{computational lawfulness}. Informally, two implementations are identified
when they realize the same abstract computation, regardless of parameterization
or microscopic event history.

Let $\mathsf{Comp}$ be a category of abstract computations (types/objects and
programs/morphisms) adequate to represent the target algorithmic structures (folds,
automata, dynamic programs, etc.).

\begin{definition}[Lawfulness quotient functor]
A \emph{lawfulness quotient} is a functor
\[
\mathcal{Q}_{\mathrm{CDL}} : \mathsf{Fld} \to \mathsf{Comp}
\]
that maps a field configuration $X$ to an abstract computational object
$\mathcal{Q}_{\mathrm{CDL}}(X)$ and maps a field update to its induced computation,
identifying any two realizations that satisfy the same pre/post-condition behavior
(or equivalently, the same algebraic laws).
\end{definition}

At this stage, parameters become explicit, and reparameterizations become
2-morphisms. One convenient formalization is via the 2-category of parametric maps.

\begin{definition}[Parametric maps (sketch)]
Let $\mathbf{Para}$ have objects $(A,P)$, where $A$ is an input space and $P$ a
parameter space. A 1-morphism $(A,P)\to(B,Q)$ is a parameterized map
$f:P\times A\to B$ (with an appropriate parameter transport to $Q$), and a
2-morphism is a reparameterization $\varphi:P\to P'$ such that the induced maps
agree up to $\varphi$.
\end{definition}

Within CDL, constraints are expressed as homomorphism conditions for algebras of
endo-functors.

\begin{definition}[Endofunctor algebras and constraint as homomorphism]
Let $F$ be an endofunctor on a category $\mathcal{C}$. An $F$-algebra is a pair
$(A,\alpha)$ with $\alpha:FA\to A$. A morphism of $F$-algebras
$f:(A,\alpha)\to(B,\beta)$ satisfies
\[
f\circ \alpha = \beta \circ Ff.
\]
In CDL, an architectural \emph{constraint} is precisely such a homomorphism
requirement.
\end{definition}

\begin{proposition}[CDL as a quotient of RSVP distinct from GDL]
The functor $\mathcal{Q}_{\mathrm{CDL}}$ identifies RSVP realizations by
computational lawfulness (algebraic semantics), not by symmetry orbits. Hence
$\mathcal{Q}_{\mathrm{CDL}}$ and $\mathcal{Q}_{\mathrm{GDL}}$ are, in general,
incomparable: they forget different structure.
\end{proposition}

\begin{proof}
$\mathcal{Q}_{\mathrm{GDL}}$ quotients by group actions, retaining invariants under
$G$. $\mathcal{Q}_{\mathrm{CDL}}$ quotients by satisfaction of algebraic laws (e.g.\
fold laws, automaton laws), retaining invariants under reparameterization and
implementation details. Since group invariants need not determine algebraic
semantics and vice versa, neither quotient refines the other in general.
\end{proof}

% --------------------------------------------------
\subsection{Summary as a commuting extraction pattern}
% --------------------------------------------------

The derivation can be summarized as the existence of a coarse-graining from
implementations to fields, followed by two distinct abstractions:
\[
\mathsf{Hist}
\;\xrightarrow{\ \mathcal{Q}_{\mathrm{RSVP}}\ }\;
\mathsf{Fld}
\;\xrightarrow{\ \mathcal{Q}_{\mathrm{GDL}}\ }\;
\mathsf{Fld}/G,
\qquad
\mathsf{Fld}
\;\xrightarrow{\ \mathcal{Q}_{\mathrm{CDL}}\ }\;
\mathsf{Comp}.
\]
In words: RSVP is the thermodynamic quotient of eventful implementations; GDL is
the symmetry quotient of RSVP; CDL is the lawfulness quotient of RSVP. The
ordering expresses the intended dependency: categorical and geometric
architectural languages are extracted from an underlying field-theoretic
description, itself extracted from implementation histories.

% ==================================================
\section{Conclusion: Constraints, Commitment, and the Boundary of Algebra}
% ==================================================

This paper has argued for a precise and limited claim: Categorical Deep Learning
provides a complete and elegant theory of \emph{lawful computation}, but lawfulness
alone is insufficient to account for commitment, refusal, history, or agency.
This is not a failure of the framework. It is the consequence of its success.

By unifying constraints and implementations through algebraic structure,
categorical deep learning explains why certain architectures generalize, why
weight tying and recursion are principled, and why invariances can be enforced
without loss of expressivity. It shows that much of modern deep learning can be
understood as the study of structure-preserving maps between parameterized
computational spaces. In doing so, it replaces heuristic design with necessity.

At the same time, the analysis has shown that everything preserved by this
unification is, by design, indifferent to history. Reparameterization symmetry,
homomorphism requirements, and algebraic quotients all function by erasing
implementation detail. Non-invertibility appears as information loss, not as
obligation. Optimization produces convergence, not settlement. What survives is
structure; what disappears is commitment.

The central contribution of this work is to make that boundary explicit and to
locate it formally. Rather than attempting to enrich categorical deep learning
with ad hoc notions of refusal or accountability, the paper reverses the usual
direction of derivation. It begins from eventful implementations, shows how a
field-theoretic description (RSVP) arises as a thermodynamic coarse-graining of
history, and then demonstrates how both Geometric Deep Learning and Categorical
Deep Learning emerge as distinct abstractions of that field-level description.
Each abstraction forgets something different: geometry forgets energetic and
entropic substance, while category theory forgets concrete event history.

Seen in this light, categorical deep learning is not an incomplete foundation,
but a correct quotient. It captures exactly what remains invariant under
lawfulness-preserving abstraction. What it cannot express—ownership, refusal,
strong negation, and worldhood—are not bugs to be patched but structures that have
already been erased by the act of abstraction itself.

The implications are twofold. First, expectations placed on purely algebraic
accounts of intelligence should be recalibrated. If computation explains how
systems act, it does not explain why they ever stop, give, or bind themselves to
the consequences of action. Second, any attempt to build accountable or
agent-like systems must either retain access to event history or reintroduce
commitment as a first-class primitive. No amount of additional algebraic
unification can recover what has been quotiented away.

In this sense, the title of the paper is literal. Constraints without commitment
are not a deficiency of contemporary deep learning theory; they are its defining
feature. Understanding that feature—where it comes from, what it explains, and
what it necessarily excludes—is a prerequisite for any serious theory of agency
beyond lawfulness.


\newpage
% ==================================================
\appendix
\section{Algebraic Constraint Versus Commitment}
% ==================================================

\begin{definition}[Algebraic Constraint]
Let $\mathcal{C}$ be a category and $F : \mathcal{C} \to \mathcal{C}$ an
endo-functor. An \emph{algebraic constraint} is the requirement that a morphism
$f : A \to B$ be a homomorphism between $F$-algebras, i.e., that the following
diagram commute:
\[
\begin{array}{ccc}
F(A) & \xrightarrow{F(f)} & F(B) \\
\downarrow & & \downarrow \\
A & \xrightarrow{f} & B
\end{array}
\]
Such constraints specify lawful computation by preserving algebraic structure.
\end{definition}

\begin{definition}[Commitment]
A \emph{commitment} is an event that restricts a system’s future admissible
morphisms in a way that cannot be undone by reparameterization or equivalence.
Formally, a commitment at time $t$ induces a strict subcategory
$\mathcal{C}_{t+1} \subset \mathcal{C}_t$ such that
\[
\mathrm{Hom}_{\mathcal{C}_{t+1}}(A,B) \subsetneq \mathrm{Hom}_{\mathcal{C}_t}(A,B)
\]
for at least one pair $(A,B)$.
\end{definition}

\begin{proposition}
Algebraic constraints do not imply commitment.
\end{proposition}

\begin{proof}
Algebraic constraints restrict morphisms to those preserving structure but do
not remove morphisms from the category itself. Reparameterizations via
two-morphisms preserve the homomorphism condition and therefore leave the set of
admissible morphisms invariant up to equivalence. No strict subcategory is
induced.
\end{proof}

This establishes that algebraic lawfulness constrains computation without
binding future action.

% ==================================================
\section{Non-Invertible Computation and the Absence of Obligation}
% ==================================================

\begin{definition}[Non-Invertible Morphism]
A morphism $f : A \to B$ in a category $\mathcal{C}$ is \emph{non-invertible} if
there exists no morphism $g : B \to A$ such that $g \circ f = \mathrm{id}_A$.
\end{definition}

\begin{definition}[Information Loss]
A computation exhibits \emph{information loss} if $|A| > |B|$ or if distinct
elements of $A$ are identified under $f$.
\end{definition}

\begin{proposition}
Information loss does not constitute obligation.
\end{proposition}

\begin{proof}
Non-invertibility characterizes mappings between representations, not relations
between agents and futures. The destruction of information under $f$ does not
restrict the existence of alternative morphisms $f' : A \to B$ nor does it
constrain future morphisms from $B$. Therefore no binding constraint on future
action is imposed.
\end{proof}

\begin{corollary}
A system may be logically irreversible while remaining temporally and normatively
reversible.
\end{corollary}

This shows that CDL’s move from groups to monoids suffices to model irreversible
computation but not irreversible commitment.

% ==================================================
\section{Parametric Maps and Symmetric Loss}
% ==================================================

\begin{definition}[Category of Parametric Maps]
The category $\mathbf{Para}$ has objects $(A,P)$, where $A$ is a space and $P$ a
parameter space. Morphisms are functions
\[
f : P \times A \to B
\]
and two-morphisms are reparameterizations $\phi : P \to Q$ such that
$f = g \circ (\phi \times \mathrm{id}_A)$.
\end{definition}

\begin{definition}[Symmetric Loss]
Loss is \emph{symmetric} if parameter updates apply uniformly across $P$ and no
subspace $P' \subset P$ is rendered inaccessible by the update rule.
\end{definition}

\begin{proposition}
$\mathbf{Para}$ admits only symmetric loss.
\end{proposition}

\begin{proof}
All admissible updates in $\mathbf{Para}$ are reparameterizations preserving the
homomorphism condition. There exists no mechanism to permanently exclude a
parameter region while preserving the categorical structure. Thus no unilateral
loss or ownership can be represented.
\end{proof}

\begin{corollary}
Game-theoretic commitment devices cannot be expressed in $\mathbf{Para}$ without
additional structure.
\end{corollary}

% ==================================================
\section{Refusal and Strong Negation}
% ==================================================

\begin{definition}[Weak Negation]
Weak negation is the exclusion of actions due to optimization criteria, such as
cost or probability, and may be reversed by changes in resources or context.
\end{definition}

\begin{definition}[Strong Negation (Refusal)]
Strong negation is the permanent exclusion of an action from the admissible
future action space, independent of optimization incentives.
\end{definition}

\begin{proposition}
Categorical constraints encode weak negation but not strong negation.
\end{proposition}

\begin{proof}
Categorical constraints restrict morphisms by lawfulness. If a morphism exists
and preserves structure, it remains admissible regardless of cost or preference.
Therefore no permanent exclusion based on refusal can be encoded.
\end{proof}

\begin{corollary}
Refusal requires historical state not representable in purely compositional
semantics.
\end{corollary}

% ==================================================
\section{Events and Temporal Asymmetry}
% ==================================================

\begin{definition}[Process]
A process is a sequence of composable morphisms in a category.
\end{definition}

\begin{definition}[Event]
An event is a transformation that induces a strict reduction in the future
admissible morphisms of the system.
\end{definition}

\begin{proposition}
Categories model processes but not events.
\end{proposition}

\begin{proof}
Categories are closed under composition. For any morphism $f : A \to B$, further
composition is always defined when types match. No intrinsic notion of terminal
exclusion exists.
\end{proof}

\begin{corollary}
Worldhood requires structure beyond categorical composition.
\end{corollary}

% ==================================================
\appendix
\section{Formal Derivation of Geometric and Categorical Deep Learning from RSVP}
% ==================================================

This appendix provides a formal account of the derivation order claimed in the
main text. In particular, it demonstrates that both Geometric Deep Learning (GDL)
and Categorical Deep Learning (CDL) arise as abstractions of an intermediate
field-theoretic description (RSVP), which itself arises as a quotient of
eventful implementations. The derivation is expressed using categories,
equivalence relations, and functors.

% --------------------------------------------------
\subsection{Eventful Implementations as a Free History Category}
% --------------------------------------------------

Let $\mathsf{Ev}$ be a set of primitive event types equipped with admissibility
rules specifying which events may follow which.

\begin{definition}[History Category]
Define the category $\mathsf{Hist}$ as follows:
\begin{itemize}
\item Objects are implementation states.
\item Morphisms $h : x \to y$ are finite admissible event sequences such that
      replaying $h$ from $x$ yields $y$.
\item Composition is concatenation of event sequences.
\item Identity morphisms are empty histories.
\end{itemize}
\end{definition}

$\mathsf{Hist}$ is a free category generated by event transitions. It records
full execution lineage and is maximally concrete.

% --------------------------------------------------
\subsection{RSVP as a Thermodynamic Quotient}
% --------------------------------------------------

Let $\mathsf{Fld}$ be a category whose objects are RSVP field configurations
$(\Phi,\mathbf{v},S)$ and whose morphisms are admissible field evolutions.

\begin{definition}[Thermodynamic Equivalence]
Define an equivalence relation $\sim_{\mathrm{RSVP}}$ on objects and morphisms of
$\mathsf{Hist}$ by
\[
x \sim_{\mathrm{RSVP}} y
\quad \Longleftrightarrow \quad
(\Phi_x,\mathbf{v}_x,S_x) = (\Phi_y,\mathbf{v}_y,S_y).
\]
\end{definition}

\begin{proposition}[RSVP as a Quotient Category]
There exists a functor
\[
\mathcal{Q}_{\mathrm{RSVP}} : \mathsf{Hist} \to \mathsf{Fld}
\]
that factors through the quotient $\mathsf{Hist}/\!\sim_{\mathrm{RSVP}}$.
\end{proposition}

\begin{proof}
By construction, $\mathcal{Q}_{\mathrm{RSVP}}$ is constant on equivalence classes
of $\sim_{\mathrm{RSVP}}$. Thus it descends uniquely to a functor on the quotient.
\end{proof}

RSVP therefore represents the maximal structure invariant under thermodynamic
coarse-graining of event histories.

% --------------------------------------------------
\subsection{Geometric Deep Learning as Symmetry Quotient of RSVP}
% --------------------------------------------------

Let $G$ be a group acting on $\mathsf{Fld}$ by field automorphisms.

\begin{definition}[Symmetry Equivalence]
Define $X \sim_G Y$ iff there exists $g \in G$ such that $Y = g \cdot X$.
\end{definition}

\begin{definition}[Symmetry Quotient Category]
Define $\mathsf{Fld}/G$ to be the orbit category whose objects are $G$-orbits
$[X]$ and whose morphisms are induced by $G$-equivariant maps.
\end{definition}

\begin{proposition}[GDL from RSVP]
Equivariant neural architectures correspond precisely to morphisms in
$\mathsf{Fld}/G$.
\end{proposition}

\begin{proof}
A map descends to the quotient iff it is invariant under the group action,
i.e.\ equivariant. This is exactly the defining condition of layers in Geometric
Deep Learning.
\end{proof}

Thus GDL is obtained by forgetting energetic and entropic detail while retaining
symmetry invariants.

% --------------------------------------------------
\subsection{Categorical Deep Learning as Lawfulness Quotient of RSVP}
% --------------------------------------------------

Let $\mathsf{Comp}$ be a category of abstract computations (types and programs).

\begin{definition}[Lawfulness Equivalence]
Define $X \sim_{\mathrm{law}} Y$ in $\mathsf{Fld}$ iff $X$ and $Y$ realize the same
algebraic semantics, i.e.\ they satisfy the same pre/post-condition behavior
under all admissible inputs.
\end{definition}

\begin{proposition}[Existence of Lawfulness Quotient]
There exists a functor
\[
\mathcal{Q}_{\mathrm{CDL}} : \mathsf{Fld} \to \mathsf{Comp}
\]
that identifies RSVP realizations up to $\sim_{\mathrm{law}}$.
\end{proposition}

\begin{proof}
Algebraic semantics define congruence classes on realizations. Identifying
realizations within each class yields the quotient category $\mathsf{Comp}$.
\end{proof}

Within this quotient, architectures are expressed as algebras of endofunctors,
and constraints are homomorphism conditions.

% --------------------------------------------------
\subsection{Incomparability of GDL and CDL}
% --------------------------------------------------

\begin{proposition}[Orthogonality of Abstractions]
The quotients $\mathsf{Fld}/G$ and $\mathsf{Comp}$ are generally incomparable:
neither refines the other.
\end{proposition}

\begin{proof}
$\mathsf{Fld}/G$ identifies realizations differing by symmetry actions but may
distinguish realizations with identical computational semantics.
Conversely, $\mathsf{Comp}$ identifies realizations with identical semantics
regardless of symmetry. Therefore neither quotient factors through the other in
general.
\end{proof}

% --------------------------------------------------
\subsection{Summary Diagram}
% --------------------------------------------------

The derivation order is summarized by the following diagram of quotient functors:
\[
\mathsf{Hist}
\;\xrightarrow{\ \mathcal{Q}_{\mathrm{RSVP}}\ }\;
\mathsf{Fld}
\;\xrightarrow{\ \mathcal{Q}_{\mathrm{GDL}}\ }\;
\mathsf{Fld}/G,
\qquad
\mathsf{Fld}
\;\xrightarrow{\ \mathcal{Q}_{\mathrm{CDL}}\ }\;
\mathsf{Comp}.
\]

Geometric Deep Learning and Categorical Deep Learning are thus extracted from RSVP
by forgetting different structure. RSVP itself is extracted from implementation
histories by forgetting event-level detail while retaining thermodynamic
invariants.

% ==================================================
\section{A Minimal Extension Admitting Commitment}
% ==================================================

This appendix defines a minimal categorical extension sufficient to represent
commitment and refusal, without abandoning the compositional semantics of
Categorical Deep Learning.

\begin{definition}[Commitment-Enriched Category]
Let $\mathcal{C}$ be a base category. A \emph{commitment-enriched category}
$\mathcal{C}^\sharp$ consists of:
\begin{itemize}
\item the objects and morphisms of $\mathcal{C}$;
\item a distinguished class of \emph{commitment morphisms}
      $c : A \to A$;
\item a preorder $\preceq$ on morphisms such that
      $f \preceq g$ denotes that $f$ is admissible whenever $g$ is.
\end{itemize}
\end{definition}

\begin{definition}[Commitment]
A morphism $c : A \to A$ is a \emph{commitment} if, for all $f : A \to B$,
\[
f \circ c \preceq f
\quad\text{but}\quad
c \circ f \not\preceq f.
\]
\end{definition}

Intuitively, a commitment restricts future admissible actions but cannot be
neutralized by precomposition.

\begin{proposition}
$\mathcal{C}^\sharp$ is not equivalent to $\mathcal{C}$.
\end{proposition}

\begin{proof}
In $\mathcal{C}$, admissibility is determined solely by type compatibility and
lawfulness. In $\mathcal{C}^\sharp$, admissibility depends on historical
commitments encoded by the preorder $\preceq$. Therefore no equivalence functor
can preserve both structures.
\end{proof}

\begin{corollary}
Commitment cannot be represented in purely compositional semantics.
\end{corollary}

This extension is minimal: no new objects are added, and composition remains
defined, but admissibility becomes history-sensitive.

% ==================================================
\section{Event-Recording Morphisms and Broken Reparameterization Symmetry}
% ==================================================

This appendix introduces a class of morphisms that record events explicitly and
thereby break the reparameterization symmetry central to parametric map
semantics.

\begin{definition}[Event-Recording Morphism]
Let $\mathbf{Para}$ be the 2-category of parametric maps. An
\emph{event-recording morphism} is a 1-morphism
\[
f : (A,P) \to (B,Q)
\]
together with a map
\[
\epsilon_f : P \to E
\]
into an event space $E$, such that distinct parameterizations induce distinct
event records.
\end{definition}

\begin{definition}[Event Sensitivity]
A morphism $f$ is \emph{event-sensitive} if there exists no nontrivial
reparameterization $\varphi : P \to P$ such that
\[
f \circ (\varphi \times \mathrm{id}_A) = f
\quad\text{and}\quad
\epsilon_f \circ \varphi = \epsilon_f.
\]
\end{definition}

\begin{proposition}
Event-sensitive morphisms break the reparameterization symmetry of
$\mathbf{Para}$.
\end{proposition}

\begin{proof}
In $\mathbf{Para}$, two implementations differing only by reparameterization are
identified by a 2-morphism. Event sensitivity ensures that such
reparameterizations alter $\epsilon_f$, preventing identification. Thus
symmetry is broken.
\end{proof}

\begin{corollary}
Weight tying and reparameterization cease to be purely structural when events
are recorded.
\end{corollary}

Event-recording morphisms provide a formal mechanism for lineage, refusal, and
irreversible commitment to appear within an otherwise algebraic framework.

% ==================================================
\section{Relation to Categorical Deep Learning}
% ==================================================

This appendix situates the present work relative to the framework developed by
Gavranović et al.\ in \emph{Categorical Deep Learning is an Algebraic Theory of All
Architectures}.

\begin{proposition}
The present analysis is conservative with respect to Categorical Deep Learning.
\end{proposition}

\begin{proof}
All constructions in the main text reduce to standard CDL semantics when
commitment morphisms and event records are omitted. In particular, the base
categories, endofunctor algebras, and parametric map semantics coincide exactly
with those of Gavranović et al.
\end{proof}

\begin{proposition}
Categorical Deep Learning is complete with respect to lawful computation.
\end{proposition}

\begin{proof}
CDL characterizes architectures as algebras satisfying homomorphism conditions.
These conditions fully determine which computations are lawful, invariant, and
replayable under reparameterization.
\end{proof}

\begin{proposition}
Categorical Deep Learning is incomplete with respect to commitment.
\end{proposition}

\begin{proof}
CDL admits all reparameterizations preserving algebraic structure. Therefore no
morphism can permanently restrict future admissible morphisms. Commitment and
refusal, as defined in Appendix~G, require admissibility relations not present in
the CDL formalism.
\end{proof}

\begin{corollary}
The present extensions are orthogonal to, not in competition with, CDL.
\end{corollary}

In summary, Categorical Deep Learning provides the correct algebraic account of
architectural lawfulness. The present work identifies a minimal extension
required to model historical binding, refusal, and accountability—phenomena
intentionally outside the scope of CDL as originally formulated.

% ==================================================
\appendix
\section{Worked Example: Carry-as-Commitment and Event-Recorded Refusal}
% ==================================================

This appendix gives a concrete toy example showing (i) how an event-recording
extension breaks reparameterization symmetry, and (ii) how a minimal commitment
operator encodes refusal as strong negation. The example is intentionally small:
a two-digit incrementer with carry. The point is not numerical performance but
the categorical distinction between \emph{lawful computation} and \emph{settled
history}.

% --------------------------------------------------
\subsection{Base computation: two-digit increment with carry}
% --------------------------------------------------

Let the state space be
\[
A \;=\; \{0,\dots,9\}\times\{0,\dots,9\},
\qquad
(a,b)\in A
\]
representing a two-digit number $10a+b$.

Define the deterministic increment-with-carry function
\[
\mathrm{inc}:A\to A
\]
by
\[
\mathrm{inc}(a,b)\;=\;
\begin{cases}
(a,\;b+1) & \text{if } b<9,\\
(a+1,\;0) & \text{if } b=9 \text{ and } a<9,\\
(0,\;0) & \text{if } (a,b)=(9,9).
\end{cases}
\]
The salient feature is the \emph{carry event}: $b=9$ forces a change in the
higher digit.

In standard categorical or programming semantics, this is simply a function. It
is lawful, total, and compositional.

% --------------------------------------------------
\subsection{A parametric implementation and its reparameterization symmetry}
% --------------------------------------------------

To illustrate the issue in parametric-map semantics, consider an implementation
that depends on a parameter $\theta$ encoding a choice of representation for the
digit pair. Let
\[
P=\{0,1\}
\]
be a trivial ``representation selector.'' For $\theta\in P$, define an
implementation
\[
f_\theta:A\to A
\]
by
\[
f_0=\mathrm{inc},
\qquad
f_1=\sigma^{-1}\circ \mathrm{inc}\circ \sigma,
\]
where $\sigma:A\to A$ is the digit-swap bijection $\sigma(a,b)=(b,a)$. Thus
$f_1$ implements the \emph{same computation} in a conjugate representation.

Package this in $\mathbf{Para}$ as the parameterized map
\[
f: P\times A \to A,
\qquad
f(\theta,x)=f_\theta(x).
\]

In ordinary $\mathbf{Para}$, $\theta$-choices related by a reparameterization are
regarded as implementation variation. In particular, the nontrivial
reparameterization $\varphi:P\to P$ given by $\varphi(0)=1,\varphi(1)=0$
witnesses that these are essentially the same family up to parameter relabeling.

This is precisely the symmetry we will now break.

% --------------------------------------------------
\subsection{Event recording: making carry an explicit trace}
% --------------------------------------------------

Let the event space be
\[
E=\{0,1\},
\]
where $1$ means ``a carry occurred'' and $0$ means ``no carry occurred.''

Define the event recorder for the underlying increment:
\[
\epsilon: A \to E,
\qquad
\epsilon(a,b)=
\begin{cases}
1 & \text{if } b=9,\\
0 & \text{otherwise.}
\end{cases}
\]
Now define an event-recording morphism as an extension of $f$:
\[
\widehat f : P\times A \to A\times E,
\qquad
\widehat f(\theta,x)=\bigl(f_\theta(x),\;\epsilon(x)\bigr).
\]

\begin{proposition}[Event recording breaks a reparameterization symmetry]
The reparameterization $\varphi:P\to P$ does not identify $\widehat f$ with
itself as a 2-morphism in the event-recording extension, unless it preserves the
event record.
\end{proposition}

\begin{proof}
In plain $\mathbf{Para}$, relabeling $\theta$ can identify two implementations if
their observable behavior agrees up to reparameterization.
In the event-recording extension, equality requires agreement on both outputs:
state and event.

But the event recorder depends on the \emph{input} $x$ (specifically on whether
$b=9$) and is therefore representation-sensitive. Under conjugation by
$\sigma$, the carry predicate changes from ``$b=9$'' to ``$a=9$'' if one insists
on recording carries in the original coordinate system. Consequently, the pair
$(f_\theta(x),\epsilon(x))$ is not invariant under the parameter swap
$\varphi$ unless the event recorder is transformed in lockstep.

Thus the reparameterization that was admissible in $\mathbf{Para}$ no longer
preserves the enriched observable, and the symmetry is broken.
\end{proof}

\noindent
Interpretation: two implementations that are computationally equivalent can be
distinguished by what they \emph{record} as having happened. This is the minimal
formal mechanism by which ``biography'' can appear.

% --------------------------------------------------
\subsection{Commitment as irreversible restriction of admissible continuations}
% --------------------------------------------------

We now add a commitment structure capturing strong negation. Let the enlarged
state include a single commitment flag:
\[
A^\sharp = A\times \{0,1\},
\]
where the second component $r$ records whether a refusal commitment is active.

Define a commitment morphism
\[
c: A^\sharp \to A^\sharp
\]
that activates refusal once a carry has occurred:
\[
c\bigl((a,b),r\bigr)
=
\begin{cases}
\bigl((a,b),1\bigr) & \text{if } b=9,\\
\bigl((a,b),r\bigr) & \text{otherwise.}
\end{cases}
\]
Intuitively: ``after you have carried, you are committed to a constraint.''

Now define an admissibility rule for subsequent actions. Let
\[
\mathsf{Step}=\{\mathrm{inc},\mathrm{reset}\}
\]
be two candidate next-actions, where $\mathrm{reset}(a,b)=(0,0)$ is a ``collapse''
operation. Lift both actions to $A^\sharp$ by acting on the $A$-component and
leaving $r$ unchanged.

Define that $\mathrm{reset}$ is \emph{forbidden} once refusal is active:
\[
\mathrm{reset}\ \text{admissible at }((a,b),r)
\quad \Longleftrightarrow \quad r=0.
\]
Equivalently, the commitment induces a strict reduction of future admissible
morphisms.

\begin{proposition}[Carry induces a commitment that forbids collapse]
There exists an initial state $x_0=((a,9),0)$ such that after composing with $c$,
the action $\mathrm{reset}$ is no longer admissible.
\end{proposition}

\begin{proof}
Let $x_0=((a,9),0)$. Since $b=9$, we have $c(x_0)=((a,9),1)$. By definition of
admissibility, $\mathrm{reset}$ is admissible iff $r=0$, but here $r=1$.
Therefore $\mathrm{reset}$ is forbidden after the commitment.
\end{proof}

\begin{corollary}[Strong negation is history-sensitive]
The forbiddenness of $\mathrm{reset}$ cannot be inferred from the instantaneous
two-digit value $(a,b)$ alone; it depends on the event-history summarized by the
commitment flag.
\end{corollary}

\noindent
Interpretation: a purely algebraic description of $\mathrm{inc}$ does not entail
that reset is forbidden after carry. The prohibition arises only once events are
recorded and commitments restrict admissible continuations.

% --------------------------------------------------
\subsection{What this example demonstrates}
% --------------------------------------------------

This toy system exhibits three separations that motivate the main text.

\begin{enumerate}
\item \textbf{Lawful computation vs.\ event biography.}
The increment map is lawful and compositional; the event record makes ``carry''
an explicit trace element not quotiented away by reparameterization.

\item \textbf{Information loss vs.\ obligation.}
The carry mechanism is a non-invertible feature of the transition, but obligation
appears only after introducing commitment that restricts future admissible moves.

\item \textbf{Weak vs.\ strong negation.}
A cost-based ``back-off'' could still allow reset if resources change.
Here, reset is forbidden regardless of utility once refusal is active.
\end{enumerate}

Thus, by adding (i) event-recording observables and (ii) a minimal commitment
operator that reduces future admissible morphisms, one obtains a precise formal
sense in which implementation histories can generate both the appearance of
refusal and the breaking of reparameterization symmetry.

\newpage
\begin{thebibliography}{99}

\bibitem{gavranovic2024cdl}
Bruno Gavranović, Paul Lessard, Andrew Dudzik, Tamara von Glehn,
João G.~M.~Araújo, and Petar Veličković.
\newblock \emph{Position: Categorical Deep Learning is an Algebraic Theory of All Architectures}.
\newblock arXiv:2402.15332, 2024.

\bibitem{velickovic2021gnn}
Petar Veličković.
\newblock \emph{Graph Neural Networks as Relational Inductive Biases}.
\newblock PhD thesis, University of Cambridge, 2021.

\bibitem{bronstein2021geometric}
Michael M.~Bronstein, Joan Bruna, Taco Cohen, and Petar Veličković.
\newblock \emph{Geometric Deep Learning: Grids, Groups, Graphs, Geodesics, and Gauges}.
\newblock IEEE Signal Processing Magazine, 38(6):18--42, 2021.

\bibitem{kiparsky2011grammaticalization}
Paul Kiparsky.
\newblock \emph{Grammaticalization as Optimization}.
\newblock In \emph{Grammaticalization and Language Change: New Reflections},
pages 14--51.
\newblock Oxford University Press, 2011.
\newblock ISBN 9780199582624.
\newblock \doi{10.1093/acprof:oso/9780199582624.003.0002}.


\bibitem{maclane1998categories}
Saunders Mac~Lane.
\newblock \emph{Categories for the Working Mathematician}.
\newblock Springer, 2nd edition, 1998.

\bibitem{awodey2010category}
Steve Awodey.
\newblock \emph{Category Theory}.
\newblock Oxford University Press, 2nd edition, 2010.

\bibitem{jacobs1999categorical}
Bart Jacobs.
\newblock \emph{Categorical Logic and Type Theory}.
\newblock Elsevier, 1999.

\bibitem{winskel1995event}
Glynn Winskel.
\newblock \emph{Models for Concurrency}.
\newblock In \emph{Handbook of Logic in Computer Science}, Volume 4.
\newblock Oxford University Press, 1995.

\bibitem{milner1989communication}
Robin Milner.
\newblock \emph{Communication and Concurrency}.
\newblock Prentice Hall, 1989.

\bibitem{brandom1994making}
Robert B.~Brandom.
\newblock \emph{Making It Explicit: Reasoning, Representing, and Discursive Commitment}.
\newblock Harvard University Press, 1994.

\bibitem{sartre1956being}
Jean-Paul Sartre.
\newblock \emph{Being and Nothingness}.
\newblock Washington Square Press, 1956.

\bibitem{heidegger1962being}
Martin Heidegger.
\newblock \emph{Being and Time}.
\newblock Harper \& Row, 1962.

\bibitem{vonNeumann1944games}
John von Neumann and Oskar Morgenstern.
\newblock \emph{Theory of Games and Economic Behavior}.
\newblock Princeton University Press, 1944.

\bibitem{schelling1960strategy}
Thomas C.~Schelling.
\newblock \emph{The Strategy of Conflict}.
\newblock Harvard University Press, 1960.

\bibitem{landauer1961irreversibility}
Rolf Landauer.
\newblock \emph{Irreversibility and Heat Generation in the Computing Process}.
\newblock IBM Journal of Research and Development, 5(3):183--191, 1961.

\bibitem{friston2010free}
Karl Friston.
\newblock \emph{The Free-Energy Principle: A Unified Brain Theory?}
\newblock Nature Reviews Neuroscience, 11(2):127--138, 2010.

\end{thebibliography}

\end{document}
